%% GuardianTx — AI4015 Agentic AI, Semester Project Phase 2
%% Conference-style paper, IEEEtran two-column.
%% Self-contained: embedded bibliography, compiles on Overleaf with a single click
%% (Menu > Compiler: pdfLaTeX). No external .bib needed.
%%
%% BEFORE SUBMISSION — set these three items (search for [SET ...]):
%%   1. Project/group code in the \title footnote and on the title page.
%%   2. The two authors' institutional emails (best-guess roll-number form left in).
%%   3. Verify every reference's author list / page numbers against the source PDF.
%%
\documentclass[conference]{IEEEtran}
\IEEEoverridecommandlockouts
\usepackage{cite}
\usepackage{amsmath,amssymb,amsfonts}
\usepackage{algorithmic}
\usepackage{graphicx}
\usepackage{textcomp}
\usepackage{xcolor}
\usepackage{booktabs}
\usepackage{array}
\usepackage{url}
\usepackage{hyperref}
\hypersetup{colorlinks=true,linkcolor=black,citecolor=black,urlcolor=blue}
\def\BibTeX{{\rm B\kern-.05em{\sc i\kern-.025em b}\kern-.08em T\kern-.1667em\lower.7ex\hbox{E}\kern-.125emX}}

\begin{document}

\title{GuardianTx: An Injection-Resilient, Risk-Calibrated Agent for\\
Pre-Signing Ethereum Transaction Defense\\
\thanks{AI4015 Agentic Artificial Intelligence --- Semester Project Phase~2.
Project code: [SET project/group code, e.g.\ G3-xx]. Base paper: Ruaro \textit{et al.},
``Approve Once, Regret Forever,'' USENIX Security 2025.}}

\author{\IEEEauthorblockN{Syed Kalbe Raza}
\IEEEauthorblockA{\textit{Dept.\ of Computer Science} \\
\textit{National University of Computer and Emerging Sciences (FAST-NUCES)}\\
Islamabad, Pakistan \\
Roll No.\ 22I-0794 \\
i220794@nu.edu.pk}
\and
\IEEEauthorblockN{Shayan Ahmed}
\IEEEauthorblockA{\textit{Dept.\ of Computer Science} \\
\textit{National University of Computer and Emerging Sciences (FAST-NUCES)}\\
Islamabad, Pakistan \\
Roll No.\ 22I-0822 \\
i220822@nu.edu.pk}}

\maketitle

\begin{abstract}
Wallet-draining attacks on Ethereum exploit a structural weakness: a user who
signs an \texttt{approve} or \texttt{transferFrom} transaction authorizes an
irreversible transfer whose true effect they cannot read. The base paper for this
work, Osprey, detects approve-controllable-transferFrom (ACT) vulnerabilities
\emph{offline} and \emph{contract-side}, offering no protection to a human at the
moment of signing. We present GuardianTx, a pre-signing investigating agent that
decodes and simulates a candidate transaction, selects which analysis tools to
invoke per transaction, and returns a calibrated \textsc{allow}/\textsc{warn}/%
\textsc{block}/\textsc{escalate} decision. The system targets three research gaps:
injection resilience (G1), risk calibration with abstention (G2), and dynamic
investigation (G3). On a 405-transaction simulated-chain corpus with fixed
development/test splits and three held-out attack families, the agent reaches
test F1 \textbf{0.966} with a \textbf{0.0\%} false-block rate and expected
calibration error (ECE) \textbf{0.081}, against F1 0.545 for the Phase-1 rule
baseline and F1 0.762 for an oracle re-implementation of the base paper. Against a
worst-case prompt-injection reasoner over a 120-case taxonomy, a decision-floor
verifier reduces attack success rate from \textbf{90.7\%} to \textbf{0.0\%} with no
loss of clean-input F1. An abstention band lifts selective accuracy from 0.973 to
0.996 at 94.3\% coverage, and dynamic tool selection matches a fixed pipeline's
accuracy (F1 0.966 vs.\ 0.940) using fewer tool calls (3.44 vs.\ 4.94). We report
all results, including partial resolutions and the system's failure modes, from
reproducible scripts.
\end{abstract}

\begin{IEEEkeywords}
LLM agents, prompt injection, blockchain security, Ethereum, transaction phishing,
confidence calibration, selective prediction, abstention.
\end{IEEEkeywords}

\section{Introduction}
Ethereum's token-approval mechanism~\cite{wood2014yellow,eip20} concentrates risk at
the instant of signing. When a user approves a contract to spend a token, that
contract may later call \texttt{transferFrom} to move the user's balance; if an
attacker controls the arguments of that call, every approver can be
drained~\cite{ruaro2025approve}. Signature-based approvals such as
\texttt{permit}~\cite{eip2612} widen the same surface.
Measurement studies place the scale of the resulting losses in the tens to hundreds
of millions of dollars, with drainer-as-a-service operations draining roughly
\$135M from 76{,}582 victims in a two-year window~\cite{he2025daas} and address
poisoning producing on the order of 270M on-chain attempts against 17M
victims~\cite{tsuchiya2025poisoning}. The common thread is that the human signer
cannot read the real effect of the payload they are about to authorize.

The base paper, Osprey~\cite{ruaro2025approve}, formalizes the ACT vulnerability
class and scans the chain for vulnerable contracts, flagging 32{,}582 candidates
and confirming 410 as exploitable. It is, however, an offline analysis that aids a
\emph{developer or auditor}; it produces no per-transaction risk score, no notion
of calibrated confidence, and no defense against adversarial manipulation of an
automated analyzer. These omissions define the opening for a user-side, pre-signing
agent.

Building such an agent as a large-language-model (LLM) system introduces its own
hazard. Agent Security Bench reports peak average attack success rates of
84.30\% against LLM agents and finds existing defenses of limited
effectiveness~\cite{zhang2025asb}; because our agent reads untrusted on-chain
strings, it is itself an injection target. This work therefore pursues three gaps
identified in our Phase-1 review:
\begin{itemize}
\item \textbf{G1 (Injection resilience).} No pre-signing defense remains robust to
prompt injection embedded in untrusted on-chain data.
\item \textbf{G2 (Calibration \& abstention).} Prior detectors report accuracy or
F1 but are neither risk-calibrated nor able to abstain, despite evidence that LLM
confidence overestimates correctness~\cite{jung2025trust}.
\item \textbf{G3 (Dynamic investigation).} No agent chooses, per transaction, which
tools to run across decoding, simulation, allowance inspection, and proxy
resolution.
\end{itemize}

\noindent\textbf{Contributions.} (1) A complete pre-signing agent with a decode
$\rightarrow$ investigate $\rightarrow$ assess $\rightarrow$ calibrate $\rightarrow$
verify pipeline, each gap addressed by a switchable component. (2) A
decision-floor output verifier that drives worst-case injection success to zero
without degrading clean-input accuracy. (3) A calibrated abstention band evaluated
with ECE, Brier score, and risk--coverage analysis. (4) A reproducible evaluation
(E1--E6) on fixed splits with held-out attack families, reported honestly including
partial results and failure modes.

\section{Related Work}
\textbf{Base paper and contract-side analysis.} Osprey~\cite{ruaro2025approve}
anchors our threat model with its ACT formalization and large-scale exploit
generation. It is contract-side and offline; GuardianTx moves defense to the user,
pre-signing. Broader surveys of Ethereum smart-contract
attacks~\cite{atzei2017survey}, empirical evaluations of automated analysis
tools~\cite{durieux2020empirical}, and the observation that vulnerable does not imply
exploited~\cite{perez2021vulnerable} frame why a signing-time, context-sensitive
decision differs from static auditing.

\textbf{Transaction-time phishing defense.} DeepTx~\cite{liu2025deeptx} is the
closest prior system: it simulates a transaction and has multiple LLMs judge intent,
but reports no injection-robustness evaluation, no calibrated confidence, and uses a
fixed feature-then-LLM pipeline. Payload-based phishing
detection~\cite{chen2025payload} achieves high accuracy with hand-crafted rules that
are brittle to unseen patterns and offer no abstention. PROXiFY~\cite{qasse2025proxify}
detects proxy contracts from bytecode (98.6\% precision), validating one of our
investigative tools. These studies motivate the approve/permit and proxy surfaces we
guard and inform our corpus construction~\cite{he2025daas,tsuchiya2025poisoning}.

\textbf{Agentic security analysis.} Multi-agent and multimodal LLM systems improve
over single models on contract vulnerability
detection~\cite{jie2025agent4vul,wei2025smartaudit}, and explanation-faithfulness
work~\cite{yu2025smartllama} shows why naive LLM rationales mislead --- directly
relevant to our verifier, which screens rationales for provenance. The agentic
paradigm itself draws on reasoning-and-acting~\cite{yao2023react} and
tool-use~\cite{schick2023toolformer,wei2022cot} methods.

\textbf{Injection and robustness.} Indirect prompt injection compromises
LLM-integrated applications through data the model
consumes~\cite{greshake2023injection,perez2022ignore}, codified as the top risk for
LLM applications~\cite{owasp2025llm}; ASB~\cite{zhang2025asb} quantifies agent-level
vulnerability. GuardianTx treats on-chain \texttt{metadata} as the adversarial
surface and hardens the decision, not merely the prompt.

\textbf{Calibration and abstention.} We use Platt scaling~\cite{niculescu2005predicting}
selected by cross-validation, measured with ECE~\cite{naeini2015obtaining} and the
Brier score~\cite{brier1950verification}, and frame the escalation band as selective
prediction~\cite{geifman2017selective}; the deferral-when-uncertain principle follows
``Trust or Escalate''~\cite{jung2025trust}.

\section{Problem Formulation}
A candidate transaction $t=(\textit{to},\textit{value},\textit{calldata},\textit{ctx})$
is submitted before signing, where \textit{ctx} includes on-chain state (balances,
allowances, code, reputation) and attacker-influenced \texttt{metadata}. The system
must output a decision $d(t)\in\{\textsc{allow},\textsc{warn},\textsc{block},
\textsc{escalate}\}$ and a risk score $p(t)\in[0,1]$. \textsc{allow}/\textsc{warn}
pass the transaction, \textsc{block} stops it, and \textsc{escalate} abstains and
defers to a human. A false block (benign $\rightarrow$ \textsc{block}) is treated as
the most damaging error because it erodes trust and is irreversible from the user's
standpoint; a false allow (malicious $\rightarrow$ pass) is the safety failure we
primarily guard against. We formalize three hypotheses:
\begin{itemize}
\item \textbf{H1 (G1).} A decision-floor verifier makes the worst-case injection
attack success rate (ASR) zero without reducing clean-input F1.
\item \textbf{H2 (G2).} A calibrated score with an abstention band lowers calibration
error and raises selective accuracy relative to a non-abstaining detector.
\item \textbf{H3 (G3).} Per-transaction tool selection matches a fixed pipeline's
accuracy with fewer tool calls and generalizes to held-out attack families.
\end{itemize}

\section{Proposed System}
\subsection{Architecture}
GuardianTx is an agent loop (Fig.~\ref{fig:arch} is described textually here for the
camera-ready): a \emph{planner} proposes the next action given the current state; a
\emph{toolbox} executes it against a simulated chain; an \emph{evidence} store
records trusted findings; and on \texttt{finish} the system runs \emph{assess}
$\rightarrow$ \emph{calibrate} $\rightarrow$ \emph{verify} to emit the decision. The
default \textsc{PolicyPlanner} is deterministic (no LLM) and data-dependent; an
\textsc{LLMPlanner} drives the same loop with a live model; a \textsc{FixedPlanner}
runs a static tool order (the G3 ablation); and a \textsc{CompliantPlanner} models a
fully hijacked reasoner for worst-case G1 evaluation.

\subsection{Environment and tools}
A deterministic simulated chain (\texttt{world.json}) stands in for a local Anvil
node so that every result is reproducible offline; it models ERC-20/721 semantics,
allowances, proxy indirection, reputation, and drainer sweeps. The agent has five
tools: calldata decoding, signature/reputation lookup, transaction simulation,
allowance inspection, and proxy resolution. All attacker-influenceable text enters
only through \texttt{metadata}; every other field is treated as trusted.

\subsection{G1 --- Injection resilience}
Untrusted strings pass a sanitizer (\texttt{guard.screen}) at a defined boundary
before any reasoning. Because a regex sanitizer generalizes poorly (Sec.~\ref{sec:g1}),
the load-bearing component is an \emph{output verifier} with three checks: a
\emph{decision floor} that prevents a reasoner from downgrading a risk level the
trusted evidence already justifies; an \emph{evidence checklist} that requires named
trusted findings for each decision; and an \emph{output-provenance} check that
withholds any reasoner rationale sharing a long verbatim span with untrusted input.

\subsection{G2 --- Calibration and abstention}
A calibrator is fit on the development split only. Platt
scaling~\cite{niculescu2005predicting} is retained solely when it beats the identity
map under development cross-validation; otherwise the identity map is kept, avoiding
test-set leakage. Two thresholds, $t_{\text{low}}=0.25$ and $t_{\text{block}}=0.475$,
define the bands; scores between them yield \textsc{escalate}. We report ECE, Brier
score, false-block rate, and a risk--coverage curve.

\subsection{G3 --- Dynamic investigation}
The \textsc{PolicyPlanner} selects the next tool from the current evidence (e.g.,
resolve a proxy (via its EIP-1967 storage slots~\cite{eip1967}) only when code
indicates indirection, inspect allowance only for approval-bearing calldata),
stopping when the decision is determined. The
\textsc{FixedPlanner} always runs the full tool sequence, isolating the value of
dynamic selection.

\section{Experimental Setup}
\subsection{Corpus and splits}
The corpus contains 405 transactions generated by \texttt{data/build\_corpus.py}
(seed 2026) over the simulated chain: a development split of 144 and a test split of
261. The test split comprises 216 \emph{seen}-family transactions and 45 from
\emph{three held-out attack families} never seen in development, isolating
generalization. The injection taxonomy has 120 cases across five categories
(authority spoofing, instruction override, tool-output laundering,
role/reasoning spoofing, multilingual/encoded) crossed with three goals (downgrade,
truncate, rationale); development and test injection templates are disjoint so that
sanitizer cues cannot be tuned on test phrasing.

\subsection{Models and settings}
The \textsc{PolicyPlanner} is deterministic and requires no model. The
\textsc{LLMPlanner} runs on Cerebras \texttt{gpt-oss-120b} (temperature 0.2,
top-$p$ 0.7), with a provider fallback order of NVIDIA NIM $\rightarrow$ Cerebras
$\rightarrow$ Groq recorded per call in each trace. As a stochastic configuration it
is repeated three times and reported as mean $\pm$ standard deviation; owing to
free-tier rate limits (approximately one minute per transaction under backoff), the
LLM planner is evaluated on a fixed, seeded, stratified sample ($n=12$).

\subsection{Baselines}
\textbf{E1} is the Phase-1 deterministic decode-then-rules pipeline. \textbf{E2} is
the base paper: Osprey cannot run in our environment (it requires contract bytecode),
so we use an \emph{oracle re-implementation} of its ACT detection with access to
ground-truth labels. Its numbers are an \emph{upper bound} for contract-side
detection, explicitly not figures reported in the paper.

\subsection{Metrics}
We report precision, recall, F1 (with 95\% confidence intervals), false-block and
false-flag rates, escalation rate, ECE, Brier score, selective accuracy and
coverage, tool-call counts, and --- for G1 --- attack success rate (ASR), the
fraction of eligible injections that achieve the attacker's goal.

\section{Results}
\subsection{E1--E3: baseline, base paper, proposed system}
Table~\ref{tab:main} compares the three systems on the test split. The rule baseline
is poorly calibrated (ECE 0.344) and blocks 25.8\% of benign transactions; the
oracle base-paper detector is precise (0.978) but misses many malicious cases
(recall 0.624) because it reasons only about contract-side ACT patterns. The
proposed agent reaches F1 0.966 with a 0.0\% false-block rate, ECE 0.081, and recall
0.993, while escalating only 5.7\% of transactions.

\begin{table}[t]
\caption{Test split ($n=261$). E2 is an oracle upper bound, not base-paper figures.}
\label{tab:main}
\centering
\small
\begin{tabular}{@{}lccccc@{}}
\toprule
System & F1 & Prec. & Rec. & FBR & ECE \\
\midrule
E1 Rule baseline        & 0.545 & 0.616 & 0.489 & 0.258 & 0.344 \\
E2 Osprey (oracle UB)   & 0.762 & 0.978 & 0.624 & 0.017 & 0.211 \\
E3 GuardianTx (policy)  & \textbf{0.966} & 0.940 & \textbf{0.993} & \textbf{0.000} & \textbf{0.081} \\
E3 GuardianTx (LLM)$^\dagger$ & 0.838 & 0.722 & 1.000 & 0.056 & 0.091 \\
\bottomrule
\end{tabular}
\\[2pt]
\raggedright\footnotesize $^\dagger$Cerebras \texttt{gpt-oss-120b}, 3 runs $\times$
stratified $n=12$, mean values; FBR and ECE standard deviations 0.096 and 0.021.
\end{table}

The deterministic policy agent is the system of record. The LLM planner, run three
times, is stable: recall was 1.000 in every run (no malicious transaction missed);
one of three runs blocked a benign transaction once, producing the nonzero
false-block mean. On the held-out families the policy agent attains F1 1.000 versus
0.500 for the rule baseline.

\subsection{E4 (G1): injection resilience}
\label{sec:g1}
Table~\ref{tab:g1} reports ASR against the worst-case compliant reasoner over 75
eligible test injections. The regex sanitizer alone is insufficient: its screening
recall is 1.00 on development payloads but only 0.20 on test payloads, so standalone
ASR falls only to 0.72. The decision-floor verifier reduces ASR to 0.0 by
construction, and the full guard holds it at 0.0, while clean-input F1 is unchanged
($\Delta=0.0$ points for the worst-case reasoner; $-1.99$ points for the live LLM,
within run-to-run noise). With a live LLM reasoner, ASR is already 0.0 because the
model declines the injected instructions --- but H1 requires robustness even when the
reasoner is fully compromised, which only the verifier provides.

\begin{table}[t]
\caption{G1 attack success rate (worst-case reasoner, 75 eligible test injections).}
\label{tab:g1}
\centering
\small
\begin{tabular}{@{}lcc@{}}
\toprule
Defense configuration & ASR & $\Delta$ clean-F1 \\
\midrule
No defense        & 0.907 & --- \\
Sanitizer only    & 0.720 & 0.0 \\
Verifier only     & \textbf{0.000} & 0.0 \\
Full guard        & \textbf{0.000} & 0.0 \\
\bottomrule
\end{tabular}
\end{table}

\subsection{E4 (G2): calibration and abstention}
On the agent, Platt scaling does not beat the identity map on development
cross-validation, so the identity map is kept and ECE is unchanged at 0.081; the
contribution of G2 is therefore the \emph{abstention band}, which raises selective
accuracy from 0.973 (no abstention, coverage 1.0) to 0.996 at 94.3\% coverage and
raises recall from 0.950 to 0.993. Applied to the weaker Phase-1 scorer, the same
calibration machinery reduces ECE from 0.344 to 0.103 and Brier from 0.353 to 0.262,
confirming the method works where miscalibration exists. We therefore rate G2
\emph{partially} addressed: the abstention mechanism is effective, but the agent's
base scores were already well-calibrated, so calibration proper has little headroom.

\subsection{E4 (G3): dynamic investigation}
The dynamic planner attains F1 0.966 on the full test split versus 0.940 for the
fixed order, using 3.44 tool calls per transaction against 4.94 --- a 30\% reduction
in tool work at higher accuracy. Both reach F1 1.000 on the held-out families, so the
generalization benefit is in efficiency rather than accuracy on these families; we
rate G3 \emph{partially} addressed.

\subsection{E5: ablation}
Table~\ref{tab:ablation} removes one component at a time. Removing the verifier
raises worst-case ASR from 0.0 to 0.72, confirming it is the load-bearing G1
component; removing the sanitizer alone leaves ASR at 0.0 because the verifier floor
still holds. Removing calibration raises F1 slightly (0.975) but eliminates
abstention (escalation 0\%) and lowers recall to 0.950, trading a safety margin for a
headline metric. Removing the simulation tool worsens ECE (0.081 $\rightarrow$ 0.121).
Disabling dynamic selection (fixed order) lowers F1 to 0.940 and raises tool calls to
4.94.

\begin{table}[t]
\caption{Ablation (test split). ASR is worst-case reasoner.}
\label{tab:ablation}
\centering
\small
\begin{tabular}{@{}lcccc@{}}
\toprule
Configuration & F1 & ECE & ASR & Tools \\
\midrule
Full system        & 0.966 & 0.081 & 0.000 & 3.44 \\
$-$ sanitizer (G1) & 0.966 & 0.081 & 0.000 & 3.44 \\
$-$ verifier (G1)  & 0.966 & 0.081 & \textbf{0.720} & 3.44 \\
$-$ calibration (G2)& 0.975 & 0.081 & 0.000 & 3.44 \\
$-$ simulation     & 0.966 & \textbf{0.121} & 0.000 & 3.28 \\
Fixed order (G3)   & \textbf{0.940} & 0.082 & 0.000 & \textbf{4.94} \\
\bottomrule
\end{tabular}
\end{table}

\subsection{E6: efficiency and reliability}
The dynamic planner averages 3.44 tool calls (max 6) versus 4.94 (max 7) for fixed
order, at sub-millisecond offline latency. Under injected transient tool failures the
agent degrades gracefully with retries: at a 30\% fault rate, F1 is
$0.948\pm0.004$ with retries versus $0.798\pm0.003$ without, and decision agreement
with the fault-free run is 96.6\% versus 53.5\%. When the LLM reasoner is fully
unavailable the agent fails safe to \textsc{escalate}.

\subsection{Gap-resolution table}
Table~\ref{tab:gap} summarizes the outcome per gap.

\begin{table*}[t]
\caption{Gap-resolution summary. Base-paper column is an oracle re-implementation (upper bound), not reported figures.}
\label{tab:gap}
\centering
\small
\begin{tabular}{@{}p{2.2cm}p{3.6cm}p{3.0cm}p{2.3cm}p{3.2cm}c@{}}
\toprule
Gap & Method / component & Metric & Base paper & Our result & Addressed? \\
\midrule
G1 Injection resilience & Sanitizer + decision-floor verifier & Worst-case ASR; clean-F1 change & No reasoning layer (N/A) & ASR 90.7\%$\rightarrow$0.0\%; F1 $\Delta$ +0.0 pts & Fully$^\ast$ \\
G2 Calibration \& abstention & CV-selected Platt + abstention band & ECE, FBR, selective accuracy & ECE 0.211, FBR 1.7\% & Selective acc.\ 0.973$\rightarrow$0.996 @ 94.3\% cov.; ECE 0.081; Phase-1 scorer ECE 0.344$\rightarrow$0.103 & Partially \\
G3 Dynamic investigation & PolicyPlanner tool selection & Held-out F1; tool calls & F1 1.000 (oracle) & F1 0.966 vs 0.940 fixed; 3.44 vs 4.94 calls; held-out F1 1.000 & Partially \\
\bottomrule
\end{tabular}
\\[2pt]
\raggedright\footnotesize $^\ast$G1 is rated fully only because the decision-floor
verifier makes downgrade attacks fail by construction even against a fully hijacked
reasoner; the regex sanitizer alone generalizes poorly (test screening recall 0.20).
\end{table*}

\section{Failure Analysis}
The final policy agent makes 10 errors on the 261 test transactions, in two
categories. Nine are false positives of type \emph{new/unverified legitimate
counterparty}: benign interactions with new decentralized applications
(\texttt{new\_dapp\_bounded}, and two \texttt{token\_transfer} cases) that trip the
\texttt{UNVERIFIED\_SPENDER} code and are escalated at risk score 0.44 --- the agent
abstains rather than blocks, so no benign transaction is hard-blocked. One is a false
negative of type \emph{attacker with clean reputation}: a \texttt{transfer\_to\_scam}
case whose counterparty had no adverse reputation signal at decision time.

\textbf{Representative success cases.} (i) A multicall proxy-drain held-out
transaction is correctly \textsc{block}ed after the agent resolves the proxy and
inspects the resulting allowance delta. (ii) A bounded approval to a verified router
is correctly \textsc{allow}ed without unnecessary escalation.
\textbf{Representative failures.} (iii) A legitimate interaction with a brand-new,
unverified dApp is \textsc{escalate}d (false positive) because reputation is absent.
(iv) A transfer to a freshly created scam address is \textsc{allow}ed (false
negative) because no reputation or simulation signal distinguishes it. Both point to
the same root cause: reliance on reputation signals that are missing for new
addresses. Future work could add behavioral priors (e.g., contract age, funding
lineage) to reduce both error types.

\section{Limitations and Threats to Validity}
Results are obtained on a \emph{simulated} chain, not live mainnet traffic; the
simulator models ERC-20/721 semantics and drainer sweeps but not the full EVM, and
trusted reputation fields are assumed correct. The base-paper comparison is an oracle
upper bound, not Osprey's own output. The LLM-planner evaluation is small ($n=12$,
three runs) due to free-tier rate limits, so its point estimates carry wide
intervals; the deterministic planner, which does not depend on a model, is the system
of record. The injection taxonomy, though disjoint across splits, is author-authored
and may not capture adaptive attackers. Finally, the decision-floor guarantee for G1
holds against reasoner compromise but not against an attacker who can forge trusted
on-chain state, which is outside our threat model.

\section{Conclusion}
GuardianTx shows that a pre-signing agent can turn a brittle rule baseline (F1 0.545,
25.8\% false blocks) into a calibrated, injection-resilient detector (F1 0.966, 0.0\%
false blocks, ECE 0.081) that surpasses an oracle upper bound of the base paper on
F1 while abstaining when uncertain. A decision-floor verifier eliminates worst-case
injection success without accuracy loss (G1, fully), an abstention band improves
selective accuracy (G2, partially), and dynamic tool selection preserves accuracy
with less work (G3, partially). Future work includes behavioral priors for unverified
counterparties, larger live-model evaluation, and validation against adaptive
injections and live chain data.

\section*{Reproducibility}
All numbers come from \texttt{results/} and are regenerated by
\texttt{python experiments/run.py all} (E1--E6, offline) and
\texttt{python experiments/run.py e3 --planner llm --runs 3 --sample 12} (LLM
planner); the corpus is built by \texttt{python data/build\_corpus.py} (seed 2026).

\begin{thebibliography}{00}
\bibitem{ruaro2025approve} N. Ruaro, F. Gritti, D. Meng, R. McLaughlin, I. Grishchenko, C. Kruegel, and G. Vigna, ``Approve Once, Regret Forever: On the Exploitation of Ethereum's Approve-TransferFrom Ecosystem,'' in \textit{Proc. 34th USENIX Security Symp.}, 2025, pp.\ 1281--1298. \url{https://www.usenix.org/conference/usenixsecurity25/presentation/ruaro}
\bibitem{liu2025deeptx} Y. Liu, X. Li, and Y. Li, ``DeepTx: Real-Time Transaction Risk Analysis via Multi-Modal Features and LLM Reasoning,'' in \textit{Proc. 40th IEEE/ACM Int. Conf. Automated Software Engineering (ASE), Tool Demonstration Track}, 2025, arXiv:2510.18438. \url{https://arxiv.org/abs/2510.18438}
\bibitem{zhang2025asb} H. Zhang \textit{et al.}, ``Agent Security Bench (ASB): Formalizing and Benchmarking Attacks and Defenses in LLM-based Agents,'' in \textit{Proc. Int. Conf. Learning Representations (ICLR)}, 2025, arXiv:2410.02644. \url{https://arxiv.org/abs/2410.02644}
\bibitem{he2025daas} B. He \textit{et al.}, ``Unmasking the Shadow Economy: A Deep Dive into Drainer-as-a-Service Phishing on Ethereum,'' in \textit{Proc. ACM Internet Measurement Conf. (IMC)}, 2025, pp.\ 430--444, doi: 10.1145/3730567.3764476. \url{https://doi.org/10.1145/3730567.3764476}
\bibitem{chen2025payload} Z. Chen, Y. Hu, B. He, D. Luo, L. Wu, and Y. Zhou, ``Dissecting Payload-based Transaction Phishing on Ethereum,'' in \textit{Proc. Network and Distributed System Security Symp. (NDSS)}, 2025, arXiv:2409.02386. \url{https://arxiv.org/abs/2409.02386}
\bibitem{tsuchiya2025poisoning} T. Tsuchiya, J.-D. Dong, K. Soska, and N. Christin, ``Blockchain Address Poisoning,'' in \textit{Proc. 34th USENIX Security Symp.}, 2025, arXiv:2501.16681. \url{https://arxiv.org/abs/2501.16681}
\bibitem{qasse2025proxify} I. Qasse, M. Hamdaqa, and B. {\TH}. J{\'o}nsson, ``PROXiFY: A Bytecode Analysis Tool for Detecting and Classifying Proxy Contracts in Ethereum,'' in \textit{Proc. ASE 2025, Tool Demonstration Track}, 2025. \url{https://conf.researchr.org/details/ase-2025/ase-2025-tool-demonstration-track/37}
\bibitem{jie2025agent4vul} W. Jie \textit{et al.}, ``Agent4Vul: Multimodal LLM Agents for Smart Contract Vulnerability Detection,'' \textit{Science China Information Sciences}, 2025, doi: 10.1007/s11432-024-4402-2. \url{https://doi.org/10.1007/s11432-024-4402-2}
\bibitem{wei2025smartaudit} Z. Wei \textit{et al.}, ``Advanced Smart Contract Vulnerability Detection via LLM-Powered Multi-Agent Systems,'' \textit{IEEE Trans. Software Engineering}, vol.\ 51, no.\ 10, pp.\ 2830--2846, 2025, doi: 10.1109/TSE.2025.3597319. \url{https://doi.org/10.1109/TSE.2025.3597319}
\bibitem{yu2025smartllama} L. Yu \textit{et al.}, ``Smart-LLaMA-DPO: Reinforced Large Language Model for Explainable Smart Contract Vulnerability Detection,'' \textit{Proc. ACM Software Engineering}, vol.\ 2, no.\ ISSTA, pp.\ 182--205, 2025, doi: 10.1145/3728878. \url{https://doi.org/10.1145/3728878}
\bibitem{jung2025trust} J. Jung, F. Brahman, and Y. Choi, ``Trust or Escalate: LLM Judges with Provable Guarantees for Human Agreement,'' in \textit{Proc. Int. Conf. Learning Representations (ICLR)}, 2025, arXiv:2407.18370. \url{https://arxiv.org/abs/2407.18370}
\bibitem{greshake2023injection} K. Greshake, S. Abdelnabi, S. Mishra, C. Endres, T. Holz, and M. Fritz, ``Not What You've Signed Up For: Compromising Real-World LLM-Integrated Applications with Indirect Prompt Injection,'' in \textit{Proc. 16th ACM Workshop on Artificial Intelligence and Security (AISec)}, 2023, pp.\ 79--90, doi: 10.1145/3605764.3623985. \url{https://doi.org/10.1145/3605764.3623985}
\bibitem{perez2022ignore} F. Perez and I. Ribeiro, ``Ignore Previous Prompt: Attack Techniques for Language Models,'' in \textit{NeurIPS ML Safety Workshop}, 2022, arXiv:2211.09527. \url{https://arxiv.org/abs/2211.09527}
\bibitem{owasp2025llm} OWASP, ``OWASP Top 10 for LLM Applications 2025,'' 2025. \url{https://genai.owasp.org/llm-top-10/}
\bibitem{yao2023react} S. Yao, J. Zhao, D. Yu, N. Du, I. Shafran, K. Narasimhan, and Y. Cao, ``ReAct: Synergizing Reasoning and Acting in Language Models,'' in \textit{Proc. Int. Conf. Learning Representations (ICLR)}, 2023, arXiv:2210.03629. \url{https://arxiv.org/abs/2210.03629}
\bibitem{schick2023toolformer} T. Schick \textit{et al.}, ``Toolformer: Language Models Can Teach Themselves to Use Tools,'' in \textit{Proc. Advances in Neural Information Processing Systems (NeurIPS)}, 2023, arXiv:2302.04761. \url{https://arxiv.org/abs/2302.04761}
\bibitem{wei2022cot} J. Wei \textit{et al.}, ``Chain-of-Thought Prompting Elicits Reasoning in Large Language Models,'' in \textit{Proc. Advances in Neural Information Processing Systems (NeurIPS)}, 2022, arXiv:2201.11903. \url{https://arxiv.org/abs/2201.11903}
\bibitem{niculescu2005predicting} A. Niculescu-Mizil and R. Caruana, ``Predicting Good Probabilities with Supervised Learning,'' in \textit{Proc. 22nd Int. Conf. Machine Learning (ICML)}, 2005, pp.\ 625--632, doi: 10.1145/1102351.1102430. \url{https://doi.org/10.1145/1102351.1102430}
\bibitem{naeini2015obtaining} M. P. Naeini, G. F. Cooper, and M. Hauskrecht, ``Obtaining Well Calibrated Probabilities Using Bayesian Binning,'' in \textit{Proc. 29th AAAI Conf. Artificial Intelligence}, 2015, pp.\ 2901--2907, doi: 10.1609/aaai.v29i1.9602. \url{https://doi.org/10.1609/aaai.v29i1.9602}
\bibitem{brier1950verification} G. W. Brier, ``Verification of Forecasts Expressed in Terms of Probability,'' \textit{Monthly Weather Review}, vol.\ 78, no.\ 1, pp.\ 1--3, 1950, doi: 10.1175/1520-0493(1950)078<0001:VOFEIT>2.0.CO;2. \url{https://doi.org/10.1175/1520-0493(1950)078\%3C0001:VOFEIT\%3E2.0.CO;2}
\bibitem{geifman2017selective} Y. Geifman and R. El-Yaniv, ``Selective Classification for Deep Neural Networks,'' in \textit{Proc. Advances in Neural Information Processing Systems (NeurIPS)}, 2017, arXiv:1705.08500. \url{https://arxiv.org/abs/1705.08500}
\bibitem{atzei2017survey} N. Atzei, M. Bartoletti, and T. Cimoli, ``A Survey of Attacks on Ethereum Smart Contracts (SoK),'' in \textit{Proc. 6th Int. Conf. Principles of Security and Trust (POST)}, 2017, pp.\ 164--186, doi: 10.1007/978-3-662-54455-6\_8. \url{https://doi.org/10.1007/978-3-662-54455-6_8}
\bibitem{perez2021vulnerable} D. Perez and B. Livshits, ``Smart Contract Vulnerabilities: Vulnerable Does Not Imply Exploited,'' in \textit{Proc. 30th USENIX Security Symp.}, 2021, pp.\ 1325--1341. \url{https://www.usenix.org/conference/usenixsecurity21/presentation/perez}
\bibitem{eip20} F. Vogelsteller and V. Buterin, ``EIP-20: Token Standard,'' Ethereum Improvement Proposals, 2015. \url{https://eips.ethereum.org/EIPS/eip-20}
\bibitem{eip1967} S. Palladino, F. Giordano, and H. Garcia, ``EIP-1967: Standard Proxy Storage Slots,'' Ethereum Improvement Proposals, 2019. \url{https://eips.ethereum.org/EIPS/eip-1967}
\bibitem{eip2612} M. Lundfall, ``EIP-2612: Permit Extension for EIP-20 Signed Approvals,'' Ethereum Improvement Proposals, 2020. \url{https://eips.ethereum.org/EIPS/eip-2612}
\bibitem{wood2014yellow} G. Wood, ``Ethereum: A Secure Decentralised Generalised Transaction Ledger,'' Ethereum Yellow Paper, 2014. \url{https://ethereum.github.io/yellowpaper/paper.pdf}
\bibitem{durieux2020empirical} T. Durieux, J. F. Ferreira, R. Abreu, and P. Cruz, ``Empirical Review of Automated Analysis Tools on 47{,}587 Ethereum Smart Contracts,'' in \textit{Proc. 42nd Int. Conf. Software Engineering (ICSE)}, 2020, pp.\ 530--541, doi: 10.1145/3377811.3380364. \url{https://doi.org/10.1145/3377811.3380364}
\end{thebibliography}

\end{document}
